\section{Du saut spectral \texorpdfstring{\(4\to2\to6\to54\)}{4 à 2 à 6 à 54} : couplage local, compression modulaire et intégration bidimensionnelle}
\label{sec:sucesor83-salto-espectral-4-2-6-54}

Pour deux représentants consécutifs \((k,k+1)\), considérons
\[
B_k=
\begin{pmatrix}
k^2&k(k+1)\\
k(k+1)&(k+1)^2
\end{pmatrix}\pmod 9.
\]

\begin{theorem}[Unicité du bloc diagonal adjacent]
\label{thm:sucesor83-app-bloque-central}
Il existe un unique bloc \(B_k\) dont les entrées diagonales coïncident. Il est
déterminé par \(k=4\) et vaut
\[
B_c=\begin{pmatrix}7&2\\2&7\end{pmatrix}.
\]
\end{theorem}

\begin{proof}
L'égalité des diagonales équivaut à
\[
k^2\equiv(k+1)^2\pmod 9,
\]
c'est-à-dire \(2k+1\equiv0\pmod 9\). Comme \(2\) est une unité de
\(\mathbb Z/9\mathbb Z\), la congruence possède l'unique solution
\(k\equiv4\pmod 9\).
\end{proof}

Par conséquent,
\[
B_c=7I+2R,
\qquad
R=\begin{pmatrix}0&1\\1&0\end{pmatrix},
\qquad R^2=I.
\]
Avec \(P_\pm=(I\pm R)/2\),
\[
B_c=9P_++5P_-.
\]

\begin{theorem}[Hiérarchie de la séparation entre somme et produit]
\label{thm:sucesor83-app-jerarquia-4-2-6-54}
Le bloc central a pour saut spectral local \(9-5=4\). Son coefficient
hors diagonale est \(2\). La compression par les \(3\) classes modulo
\(3\) produit la différence agrégée \(51-45=6\), et le déploiement sur les
\(9\) positions produit le défaut global
\[
9\cdot6=459-405=54.
\]
Les entiers \(4,2,6,54\) appartiennent, respectivement, au saut spectral,
au couplage local, à la compression modulaire et à l'intégration
bidimensionnelle.
\end{theorem}

\begin{proof}
La décomposition \(B_c=7I+2R=9P_++5P_-\) donne le saut et le coefficient
local. Les matrices agrégées \(M_\Pi,M_\Sigma\) satisfont
\(\sum M_\Pi=51\) et \(\sum M_\Sigma=45\). Chaque entrée agrégée représente
\(9\) positions, si bien que le déploiement multiplie la différence par
\(9\) et récupère les totaux \(459\) et \(405\).
\end{proof}

Leurs concaténations ordonnées produisent les blocs \(72\) et \(27\) :
\[
72+27=99,
\qquad
72-27=45=\det B_c.
\]
La concaténation des entrées diagonales et antidiagonales produit,
respectivement, \(77\) et \(22\). Par conséquent,
\[
72+27=77+22=99,
\qquad
77-22=55=54+1.
\]
Ces égalités expriment deux partitions exactes du même bloc ; elles
n'identifient pas à elles seules une constante externe.
